\documentclass[11pt]{article}
\usepackage[a4paper,margin=1in]{geometry}
\usepackage{fontspec}
\usepackage[boldfont=false]{xeCJK}
\setCJKmainfont{FandolSong-Regular.otf}
\usepackage{amsmath}
\usepackage{amssymb}
\usepackage{array}
\usepackage{booktabs}
\usepackage{graphicx}
\usepackage{adjustbox}
\usepackage{enumitem}
\setlist[itemize]{topsep=2pt,partopsep=0pt,parsep=0pt,itemsep=2pt,leftmargin=1.4em}
\setlist[enumerate]{topsep=2pt,partopsep=0pt,parsep=0pt,itemsep=2pt,leftmargin=1.6em}
\usepackage{parskip}
\usepackage{fvextra}
\fvset{breaklines=true,breakanywhere=true,fontsize=\small}
\setlength{\parindent}{0pt}

\begin{document}

\begin{center}
  {\LARGE\bfseries Coordination and response}\\[0.35em]
  {\large Igcse Biology}
\end{center}
\vspace{0.6em}

\section*{The nervous system}

Your body must \textbf{coordinate} 协调 all its parts and react to changes. The \textbf{nervous system} 神经系统 does this using fast \textbf{electrical impulses} 电脉冲 that travel along \textbf{neurones} 神经元 (nerve cells). It has two parts:

\begin{itemize}
\item the \textbf{central nervous system} 中枢神经系统 (CNS) — the \textbf{brain} 大脑 and the \textbf{spinal cord} 脊髓.

\item the \textbf{peripheral nervous system} 周围神经系统 (PNS) — all the nerves outside the brain and spinal cord.

\end{itemize}

The nervous system coordinates and \textbf{regulates} 调节 the body's functions.

\section*{Neurones and the reflex arc}

There are three kinds of neurone:

\begin{itemize}
\item \textbf{sensory neurones} 感觉神经元 — carry impulses from receptors to the CNS.

\item \textbf{relay neurones} 中间神经元 — pass impulses on inside the CNS.

\item \textbf{motor neurones} 运动神经元 — carry impulses from the CNS to effectors.

\end{itemize}

A \textbf{reflex} 反射 action is a fast, automatic \textbf{response} 反应 to a \textbf{stimulus} 刺激 (a change). It follows a fixed path, the \textbf{reflex arc} 反射弧:

\textbf{receptor} 感受器 → sensory neurone → relay neurone → motor neurone → \textbf{effector} 效应器 (a \textbf{muscle} 肌肉 or \textbf{gland} 腺体)

\section*{Synapses (Supplement)}

A \textbf{synapse} 突触 is a junction (a tiny gap) between two neurones. At the end of one neurone are \textbf{vesicles} 囊泡 holding \textbf{neurotransmitter} 神经递质 molecules. Across the \textbf{synaptic gap} 突触间隙, the next neurone has \textbf{receptor proteins} 受体蛋白.

Events at a synapse:

\begin{enumerate}
\item an impulse arrives and makes the vesicles release neurotransmitter into the synaptic gap.

\item the neurotransmitter molecules \textbf{diffuse} 扩散 across the gap.

\item they bind to the receptor proteins on the next neurone.

\item this starts a new impulse in the next neurone.

\end{enumerate}

Synapses make impulses travel in \textbf{one direction only}.

\section*{Sense organs and the eye}

\textbf{Sense organs} 感觉器官 are groups of receptor cells that respond to one kind of stimulus — light, sound, touch, temperature or chemicals. The eye responds to light.

\subsection*{Parts of the eye}

\begin{center}
\adjustbox{max width=\linewidth}{%
\begin{tabular}{ll}
\toprule
\textbf{Part} & \textbf{Function} \\
\midrule
\textbf{cornea} 角膜 & \textbf{refracts} 折射 (bends) the light as it enters \\
\textbf{iris} 虹膜 & controls how much light enters the \textbf{pupil} 瞳孔 \\
\textbf{lens} 晶状体 & focuses the light onto the retina \\
\textbf{retina} 视网膜 & contains light receptors, some sensitive to different colours \\
\textbf{optic nerve} 视神经 & carries the impulses to the brain \\
\textbf{blind spot} 盲点 & where the optic nerve leaves the eye; it has no receptors \\
\bottomrule
\end{tabular}}
\end{center}

\subsection*{The pupil reflex}

In bright light the pupil gets \textbf{smaller}, to protect the retina. In dim light it gets \textbf{wider}, to let more light in. \textbf{(Supplement)} This is done by two sets of \textbf{antagonistic} 拮抗 muscles in the iris: the \textbf{circular muscles} 环肌 and the \textbf{radial muscles} 辐射肌. In bright light the circular muscles contract and the pupil narrows; in dim light the radial muscles contract and the pupil widens.

\subsection*{Accommodation (Supplement)}

\textbf{Accommodation} 视觉调节 is changing the shape of the lens to focus on near or far objects, using the \textbf{ciliary muscles} 睫状肌 and the \textbf{suspensory ligaments} 悬韧带:

\begin{itemize}
\item near object: the ciliary muscles contract, the ligaments slacken, and the lens becomes fatter — more refraction.

\item far object: the ciliary muscles relax, the ligaments pull tight, and the lens becomes thinner — less refraction.

\end{itemize}

\subsection*{Rods and cones (Supplement)}

The retina has two kinds of light receptor:

\begin{itemize}
\item \textbf{rods} 视杆细胞 — very sensitive, good for \textbf{vision} 视觉 in dim light (night vision); they do not detect colour.

\item \textbf{cones} 视锥细胞 — three kinds, each absorbing a different colour, giving colour vision; they work best in bright light.

\end{itemize}

Cones are packed most densely at the \textbf{fovea} 中央凹, the part of the retina that gives the sharpest image.

\section*{Hormones}

A \textbf{hormone} 激素 is a chemical, made by a gland and carried in the blood, that changes the activity of one or more \textbf{target organs} 靶器官. Hormones are made by \textbf{endocrine glands} 内分泌腺, which release them straight into the blood.

\begin{center}
\adjustbox{max width=\linewidth}{%
\begin{tabular}{lll}
\toprule
\textbf{Gland} & \textbf{Hormone} & \textbf{Main effect} \\
\midrule
\textbf{adrenal glands} 肾上腺 & \textbf{adrenaline} 肾上腺素 & prepares the body for action \\
\textbf{pancreas} 胰腺 & \textbf{insulin} 胰岛素 (and \textbf{glucagon} 胰高血糖素) & control blood \textbf{glucose} 葡萄糖 \\
\textbf{testes} 睾丸 & \textbf{testosterone} 睾酮 & male development \\
\textbf{ovaries} 卵巢 & \textbf{oestrogen} 雌激素 & female development \\
\bottomrule
\end{tabular}}
\end{center}

Adrenaline is released in 'fight or flight' situations, such as fear or danger. It increases the breathing rate, the \textbf{heart rate} 心率, the blood glucose \textbf{concentration} 浓度 and the pupil diameter, getting the body ready to act.

\section*{Nervous control vs hormonal control}

\begin{center}
\adjustbox{max width=\linewidth}{%
\begin{tabular}{lll}
\toprule
\textbf{} & \textbf{Nervous control} & \textbf{Hormonal control} \\
\midrule
speed & very fast & slower \\
\textbf{duration} 持续时间 & short-lived & longer-lasting \\
how it travels & electrical impulses along neurones & chemicals in the blood \\
\bottomrule
\end{tabular}}
\end{center}

\section*{Homeostasis}

\textbf{Homeostasis} 稳态 is keeping a constant internal \textbf{environment} 环境 inside the body — for example a steady temperature and a steady blood glucose level.

\textbf{(Supplement)} Homeostasis works by \textbf{negative feedback} 负反馈: when a value moves away from its \textbf{set point} 设定点, the body acts to bring it back. The change itself triggers the correction.

\subsection*{Controlling blood glucose}

The \textbf{liver} 肝脏 and the pancreas keep blood glucose steady:

\begin{itemize}
\item \textbf{insulin} lowers blood glucose after a meal — the liver stores the glucose as \textbf{glycogen} 糖原.

\item \textbf{(Supplement) glucagon} raises blood glucose between meals — the liver releases glucose.

\end{itemize}

In \textbf{Type 1 diabetes} 糖尿病 the pancreas cannot make insulin, so blood glucose rises too high. It is treated with insulin injections and a carefully planned diet.

\subsection*{Controlling body temperature (Supplement)}

The brain detects the body \textbf{temperature} 温度 and keeps it steady. The skin helps:

\begin{center}
\adjustbox{max width=\linewidth}{%
\begin{tabular}{ll}
\toprule
\textbf{Structure} & \textbf{Role} \\
\midrule
\textbf{sweat glands} 汗腺 & make sweat; as it evaporates it cools the skin (\textbf{sweating} 出汗) \\
\textbf{hair erector muscles} 立毛肌 & raise the hairs to trap a layer of air \\
\textbf{blood vessels} 血管 & change how much blood flows near the skin surface \\
\textbf{fatty tissue} 脂肪组织 & gives \textbf{insulation} 隔热 to reduce heat loss \\
\bottomrule
\end{tabular}}
\end{center}

When you are \textbf{too hot}: you sweat, and the \textbf{arterioles} 小动脉 widen (\textbf{vasodilation} 血管舒张) so more blood flows to the skin \textbf{capillaries} 毛细血管 and loses heat.

When you are \textbf{too cold}: you \textbf{shiver} 颤抖 (the muscles make heat), and the arterioles narrow (\textbf{vasoconstriction} 血管收缩) so less blood reaches the surface, keeping heat in.

\section*{Tropisms}

A \textbf{tropism} 向性 is a plant's growth response to a direction:

\begin{itemize}
\item \textbf{gravitropism} 向地性 — growth in response to \textbf{gravity} 重力 (roots grow down, shoots grow up).

\item \textbf{phototropism} 向光性 — growth in response to the direction of light (shoots grow towards the light).

\end{itemize}

\textbf{(Supplement)} These responses are controlled by a chemical called \textbf{auxin} 生长素:

\begin{itemize}
\item auxin is made in the \textbf{shoot} 嫩枝 tip.

\item it diffuses down through the plant.

\item light or gravity makes the auxin spread unequally to one side.

\item auxin makes cells \textbf{elongate} 伸长 (grow longer). The side with more auxin grows faster, so the shoot bends.

\end{itemize}

\section*{Exam tips}

\begin{itemize}
\item Nervous system: CNS (brain + spinal cord) and PNS (the nerves). Impulses travel along neurones.

\item Learn the reflex arc in order: receptor → sensory → relay → motor → effector.

\item At a synapse, neurotransmitter diffuses across the gap — this makes impulses go one way only.

\item Eye: cornea and lens refract light; the lens changes shape (accommodation); rods for dim light, cones for colour.

\item Hormones are slower but longer-lasting than nerves. Insulin lowers blood glucose; glucagon raises it.

\item Too hot → sweating and vasodilation; too cold → shivering and vasoconstriction.

\item Auxin from the shoot tip controls phototropism and gravitropism by making cells elongate.

\end{itemize}


\end{document}
